\section{Every nonmaximal quadratic order at an odd prime}
\label{sec:odd-quadratic-orders}

A second coefficient boundary strengthens the parity invariant when the
coefficient ring is a quadratic order.  The useful extra operator is the
canonical homotopy for multiplication by the order of a cyclic group.
It forces certain primary lengths to be even, so that their halves
provide a new invariant modulo two.

\begin{theorem}\label{thm:odd-quadratic-orders}
Let $p$ be an odd prime, let $O=\Z_{(p)}$ or $\Z_p$, and let $d\in p^2O$.
Put $A=O[t]/(t^2-d)$.  Then the degree-$p$ kernel of generic
specialization contains a direct summand isomorphic to
\[
 \bigoplus_{i\ge1}\mathbb Z/2.
\]
Every displayed generator is represented by the difference between an
induced explicit lattice and a multiple of the regular lattice.
In particular $\Phi_A$ is not surjective.
If $O=\Z_p$ and $d$ is nonzero and nonsquare in $\Q_p$, the same direct
summand lies in the kernel of extension to the normalization of $A$.
\end{theorem}

\begin{corollary}\label{cor:odd-quadratic-orders}
For every odd prime $p$, every nonmaximal $\mathbb Z_p$-order in a
quadratic extension of $\mathbb Q_p$ fails polynomial generation.
Its degree-$p$ normalization kernel contains a split countable direct
sum of copies of $\mathbb Z/2$.
\end{corollary}

We first construct the invariant, keeping all coefficient and
exact-sequence compatibilities explicit.  Write $b=d/p^2$.

\subsection{Three odd operators on Tate cohomology}

Let $L$ be an $A$-free $C_p$-lattice and set $V=L/(p,t)L$.
The assumption $d\in p^2O$ gives two natural exact sequences
\begin{align}
 0\longrightarrow V\xrightarrow{p}L/(t,p^2)L
       \longrightarrow V\longrightarrow0,\label{eq:order-beta-seq}\\
 0\longrightarrow V\xrightarrow{t}L/pL
       \longrightarrow V\longrightarrow0.\label{eq:order-delta-seq}
\end{align}
Indeed, the middle coefficient rings are respectively
$O/p^2O$ and $\mathbb F_p[t]/t^2$, and $L$ is free.
Let $\beta$ and $\delta$ be their degree-one Tate connecting maps.
All Tate groups in this subsection are finite-dimensional over
$\mathbb F_p$.

Choose an $A$-basis and write the differentials of the complete
period-two Tate complex as
\[
 d_i=X_i+tY_i,
\]
where $X_i,Y_i$ have entries in $O$.  The differentials alternate
between $T-1$ and $N(T)=1+T+\cdots+T^{p-1}$.  The identity
$d_{i+1}d_i=0$ gives
\begin{equation}\label{eq:order-XY-relations}
 X_{i+1}X_i=-p^2bY_{i+1}Y_i,
 \qquad X_{i+1}Y_i+Y_{i+1}X_i=0.
\end{equation}
For a cycle represented by an integral column $z$ with $X_i z=pw$,
the two connecting maps are
\[
 \beta[z]=[w],\qquad \delta[z]=[Y_i z].
\]
They are the connecting maps of
\eqref{eq:order-beta-seq}--\eqref{eq:order-delta-seq}, so the formulas
are independent of lifts and bases.  For example, replacing $z$ by
$z+ph$ changes $w$ by the boundary $X_i h$; changes by coboundaries
are handled by \eqref{eq:order-XY-relations}.
Those same identities give
\begin{equation}\label{eq:order-beta-delta}
 \beta^2=-\overline b\,\delta^2,
 \qquad \beta\delta+\delta\beta=0.
\end{equation}
In fact $X_{i+1}w=-pbY_{i+1}Y_i z$, while
$X_{i+1}Y_i z=-pY_{i+1}w$, and division by $p$ gives the two formulas.

Both operators anticommute with the Tate connecting map $\partial$
of an admissible sequence of $A$-lattices.  Here is a direct check of
the signs.  Split the sequence over $A$ and write the upper block of
$d_i$ as $Z_i+tW_i$.  The upper blocks of $d_{i+1}d_i=0$ are
\begin{align*}
 X'_{i+1}Z_i+Z_{i+1}X''_i
   &=-p^2b(Y'_{i+1}W_i+W_{i+1}Y''_i),\\
 X'_{i+1}W_i+Y'_{i+1}Z_i+Z_{i+1}Y''_i+W_{i+1}X''_i&=0.
\end{align*}
On $X''_i z=pw$, the first formula, divided by $p$ and reduced
modulo $p$, gives $\beta'\partial=-\partial\beta''$.
The second says that
$\delta'\partial+\partial\delta''$ is null-homotopic, with homotopy
given by the upper blocks $W_i$.

For the third operator put
\[
 Q(z)=\frac{N(z)-p}{z-1}\in\mathbb Z[z].
\]
Define a degree-minus-one homotopy on the complete complex by
\[
 h_i=I\quad(i\text{ even}),\qquad
 h_i=-Q(T)\quad(i\text{ odd}).
\]
The polynomial identity $N(T)-(T-1)Q(T)=pI$ gives
\begin{equation}\label{eq:order-p-homotopy}
 d_{i-1}h_i+h_{i+1}d_i=pI.
\end{equation}
After reduction to $V$, this is an odd cohomology operator, using
periodicity to identify degrees minus one and plus one.  It is natural
in $V$, since its matrices are polynomials in the group action.
The group generator acts as the identity on Tate cohomology, whence
\begin{equation}\label{eq:order-h-square}
 h^2=-Q(1)=-p(p-1)/2=0.
\end{equation}
This uses the assumption that $p$ is odd.

Its relations with the two coefficient boundaries are
\begin{equation}\label{eq:order-h-boundaries}
 \beta h+h\beta=I,\qquad \delta h+h\delta=0.
\end{equation}
For the first, use \eqref{eq:order-p-homotopy} modulo $p^2$ in the
coefficient complex defining $\beta$.  If $dz=pw$, then
$d(hz)=p(z-hw)$, giving the stated formula after division by $p$.
The constant terms of the full homotopy identity differ from this
calculation only by multiples of $d=p^2b$, which vanish after that
division and reduction.
For the second, expand the full homotopy as $h_i=H_i+tK_i$.
The coefficient of $t$ in \eqref{eq:order-p-homotopy} is
\[
 X_{i-1}K_i+Y_{i-1}H_i+H_{i+1}Y_i+K_{i+1}X_i=0.
\]
Thus $\delta h+h\delta$ is null-homotopic, with the indicated $K_i$.

Finally $h$ anticommutes with $\partial$ for an admissible sequence.
For completeness, reduce the sequence modulo $(p,t)$ and write
its differential and homotopy in upper triangular blocks.  If $R_i$
is the upper block of $h_i$, the upper block of $dh+hd=0$ is
\[
 X'_{i-1}R_i+Z_{i-1}h''_i+h'_{i+1}Z_i+R_{i+1}X''_i=0.
\]
It gives $h'\partial+\partial h''=0$ on cohomology.

\subsection{An additive half-length invariant}

Put $u=\delta^2$ and take the invertible Fitting part of each Tate
group.  Fitting extraction is exact on finite-dimensional modules.
On this part $\delta$ is invertible.  Define even operators
\[
 C=\delta^{-1}\beta,\qquad J=\delta h.
\]
They commute with $u$ and with every map in the Tate long exact
sequence: each of their two odd factors anticommutes with connecting
maps.  The preceding identities give
\begin{equation}\label{eq:order-CJ}
 C^2=\overline b,\qquad J^2=0,\qquad CJ+JC=-I.
\end{equation}
In particular $\ker J=\operatorname{im}J$.  If $Jv=0$, then
$v=-JCv$; the reverse inclusion follows from $J^2=0$.
This holds on every $u$-primary submodule and every connecting-image
submodule, because all these submodules are stable under $C$ and $J$.

For a monic irreducible $f\in\mathbb F_p[u]$, $f\ne u$, let
$H^0_f$ denote the $f$-primary part.  The exact sequence
\[
 0\longrightarrow\operatorname{im}J\longrightarrow H^0_f
   \xrightarrow{J}\operatorname{im}J\longrightarrow0
\]
shows that its $\mathbb F_p[u]$ composition length is even.  Hence
\begin{equation}\label{eq:order-half-length}
 \mu_f(L)=\frac12\operatorname{length}_{\mathbb F_p[u]}H^0_f
                    \pmod2
\end{equation}
is defined.  It is additive.  Indeed, exact primary extraction gives a
six-term exact cycle.  Its two connecting-image modules have equal
ordinary composition length, because $\delta$ identifies them, and
each length is even by \eqref{eq:order-CJ}.  The even-parity length
defect is the negative of the sum of these lengths.  Dividing by two
therefore gives an even integer.

Every lattice extended from $O$ has $\delta=0$, so all these invariants
vanish on such lattices and, in particular, on Young permutation
lattices.  Only finitely many primary factors occur for any one lattice.

\subsection{The lattice and its exact order}

Choose a primitive element $\eta$ of the residue field
$k_s=\F_{p^s}$, and lift its monic irreducible polynomial over
$\F_p$ to a monic polynomial $F\in O[X]$.  Put $O_s=O[X]/(F)$.
This is a finite free local DVR over $O$, of degree $s$ and residue
field $k_s$.  Indeed, its reduction modulo $p$ is a field and
integrality forces every maximal ideal to contract to $pO$, so its
unique maximal ideal is $pO_s$.  The monic polynomial $F$ is
irreducible over $\operatorname{Frac}(O)$: a monic factorization
would lie in $O[X]$, since $O$ is integrally closed, and would
contradict irreducibility after reduction.  Thus $O_s$ is a domain.
Every nonzero element is a power of $p$ times a unit, by taking
the minimum valuation of its coefficients in the monic basis.
This proves the DVR assertion without assuming completeness.

We choose units $\alpha,\kappa\in O_s$ with
$\overline\kappa^{-1}=\eta$ and
\begin{equation}\label{eq:order-alpha}
 \alpha^p=1-d/\kappa.
\end{equation}
If $d=0$, take $\alpha=1$ and any lift of $\eta^{-1}$ for $\kappa$.
Otherwise write $d=p^m b_0$ with $m\ge2$ and $b_0\in O^\times$.
Choose a unit $y\in O_s$ lifting $-\eta\overline b_0$, and put
\[
 \alpha=1+p^{m-1}y,\qquad
 \kappa=\frac{d}{1-\alpha^p}.
\]
The quotient belongs to $O_s^\times$.  In fact
\[
 1-\alpha^p=-p^m yU,\qquad
 U=1+\sum_{j=2}^{p}\binom pj p^{j(m-1)-m}y^{j-1}
       \in1+pO_s.
\]
For $2\le j<p$ the summand coefficient has valuation
$(j-1)(m-1)\ge1$; for $j=p$ its valuation is
$p(m-1)-m\ge1$, since $p$ is odd.  Therefore
$\kappa=-b_0/(yU)$ is a unit,
$\overline\kappa^{-1}=-\overline y/\overline b_0=\eta$,
and \eqref{eq:order-alpha} holds exactly.

Let $\Lambda_s=O_s[z]/(z^p-1)$ and put
\begin{equation}\label{eq:order-D-matrix}
 q(z)=\frac{z^p-\alpha^p}{z-\alpha},\qquad
 D=\begin{pmatrix}0&z-\alpha\\\kappa q(z)&0\end{pmatrix}.
\end{equation}
On $E=\Lambda_s^2$, let the generator of $C_p$ act by $z$ and let
$t$ act by $D$.  Then
\[
 D^2=\kappa(z^p-\alpha^p)I=dI,
\]
so this is a module over $A_s=O_s[t]/(t^2-d)$.
It is free of rank $p$ over $A_s$.  Its quotient modulo $(p,t)$ is
\[
 k_s[z]/(z-1)\ \oplus\ k_s[z]/(z-1)^{p-1},
\]
which has dimension $p$.  Nakayama gives a surjection $A_s^p\to E$.
Both modules are free of rank $2p$ over $O_s$, so its kernel is zero.

This also gives an explicit $p$ by $p$ matrix over $A_s$.
Let $C_\alpha$ be the companion matrix of $q$, let $v$ be its first
coordinate column, and let $w$ be the last coordinate row.  In the
basis consisting of the first $\Lambda_s^2$ generator and
$1,z,\ldots,z^{p-2}$ times the second, the group generator is
\begin{equation}\label{eq:order-T-matrix}
 T=\begin{pmatrix}
       \alpha&\kappa^{-1}t w\\
       t v&C_\alpha
    \end{pmatrix}.
\end{equation}
Its characteristic polynomial is
$(z-\alpha)q(z)-\kappa^{-1}t^2=z^p-1$, which verifies $T^p=I$ directly.

For the coefficient boundary of $E$, both Tate parities of
$V=E/(p,t)E$ have dimension two over $k_s$.  In even degree take
$v_1,(z-1)^{p-2}v_2$, and in odd degree take $v_1,v_2$.
Modulo $p$ the off-diagonal entries of $D$ are
$z-1$ and $\overline\kappa(z-1)^{p-1}$.  The difference and norm
calculations therefore give, in both directions,
\[
 \delta=\begin{pmatrix}0&\overline\kappa^{-1}\\1&0\end{pmatrix},
 \qquad u=\overline\kappa^{-1}I.
\]

Let $\sigma$ send $t$ to $-t$.  The quadratic coefficient sequence is
\begin{equation}\label{eq:order-sign-sequence}
 0\longrightarrow E^\sigma\longrightarrow
 A_s\otimes_{O_s}\operatorname{Res}_{O_s}E
 \longrightarrow E\longrightarrow0.
\end{equation}
The final map is evaluation, and the kernel map is
$v\mapsto t\otimes v-1\otimes Dv$.
In the two coefficient coordinates its image is $(-Dv,v)$.
Multiplication by $t$ on the kernel corresponds to $-D$ on the source,
because $D^2=dI$; the sequence is split over $A_s$, since its quotient
is free.  The matrix $\operatorname{diag}(1,-1)$ anticommutes with
$D$ and commutes with $z$, so $E^\sigma\simeq E$.
Its underlying $O_s[C_p]$-module is free of rank two, whence
\begin{equation}\label{eq:order-exact-two}
 2\bigl([E]-[A_s[C_p]]\bigr)=0.
\end{equation}

\subsection{Induction and the direct summand}

Choose distinct positive integers $s_i$ tending to infinity.  In
$k_i=\mathbb F_{p^{s_i}}$ choose a primitive multiplicative element
$\eta_i$, and choose $\alpha_i,\kappa_i$ by the preceding construction.
Let $f_i$ be the
minimal polynomial of $\eta_i$ over $\mathbb F_p$.
For the corresponding lattice, the even Tate group is $k_i^2$ with
$u=\eta_i$, so its $f_i$-primary composition length is two and
$\mu_{f_i}=1$.

These scalar orbits are free.  If
$j\eta_i=\eta_i^{p^l}$ with $j\in\mathbb F_p^\times$ and
$0<l<s_i$, then
\[
 p^{s_i}-1\mid(p-1)(p^l-1),
\]
although the positive integer on the right is smaller.  If $l=0$,
the equality forces $j=1$.  Distinct degrees give disjoint orbits.

Restrict the constructed lattice from $A_{s_i}$ to $A$, call it $L_i$,
and put
\[
 x_i=[\operatorname{Ind}_{C_p}^{S_p}L_i]-s_i[A[S_p]].
\]
Equation~\eqref{eq:order-exact-two} gives $2x_i=0$.
The generator-change comparison for cyclic Tate complexes scales
$u=\delta^2$ by the corresponding element of $\mathbb F_p^\times$.
In the Mackey restriction, the normalizer supplies each such twist
once; every other intersection is trivial and supplies a free
$A[C_p]$-lattice.  Consequently
\[
 (\mu_{f_j}\circ\operatorname{Res}_{C_p}^{S_p})(x_i)=\delta_{ij}.
\]
The finite-support observation following
\eqref{eq:order-half-length} makes the direct sum of these invariants
a retraction of the map sending the standard basis of
$\bigoplus_i\mathbb Z/2$ to the $x_i$.  This proves both exact order
two and the direct-summand assertion.  Their nonzero invariants,
which vanish on every Young lattice, prove non-surjectivity.

For generic vanishing, write $K=\operatorname{Frac}(O)$.
Then $A[1/p]$ is a field, a product of two fields,
or the dual-number algebra over $K$, according as $d$ is
nonsquare, nonzero square, or zero.  The characteristic-zero
coefficient result gives torsion-free lattice Grothendieck groups
in the local cases.  In the product case representations decompose
componentwise and the classes of coefficient-free objects form a
subgroup of the product of the two characteristic-zero representation
groups, so there is again no torsion.  Thus all $x_i$ vanish generically.
When $O=\Z_p$ and $A$ is a domain, its normalization is a DVR and its
symmetric-group lattice Grothendieck group is torsion-free by
Theorem~\ref{thm:regular-local}; normalization extension therefore also
kills the $x_i$.  This completes the proof of
Theorem~\ref{thm:odd-quadratic-orders}.

\begin{proof}[Proof of Corollary~\ref{cor:odd-quadratic-orders}]
Let $O_K$ be the ring of integers of a quadratic extension of
$\mathbb Q_p$.  Since $2$ is a unit, the trace projection
$\frac12\operatorname{Tr}_{K/\mathbb Q_p}$ splits
$O_K=O\oplus O\theta$, where $\theta$ has trace zero.
Thus $\theta^2=c\in O$.  The valuation of $c$ is zero or one:
if $p^2\mid c$, the element $\theta/p$ would be integral and hence
belong to $O_K$, contradicting this integral basis.
If $c$ is a unit, its residue is nonsquare, since otherwise Hensel's
lemma would make $c$ a square in $\mathbb Q_p$.

For any order $B\subset O_K$, the same trace projection shows that
its trace-zero part is a submodule $p^rO\theta$, for some $r\ge0$.
Hence $B=O\oplus p^rO\theta$.  The nonmaximal orders have $r\ge1$
and are isomorphic to
\[
 O[t]/(t^2-p^{2r}c),
\]
whose constant term is divisible by $p^2$.  Apply the theorem.
\end{proof}

\begin{corollary}\label{cor:localized-odd-quadratic-orders}
Let $p$ be odd.  Every nonmaximal finite $\Z_{(p)}$-order in a
quadratic number field is local and fails polynomial generation.
Its degree-$p$ generic-specialization kernel contains a split
countably infinite direct sum of copies of $\Z/2$.
\end{corollary}

\begin{proof}
Put $O=\Z_{(p)}$, and let $D$ be its integral closure in the quadratic
field.  This is finite free of rank two over the DVR $O$.
The trace projection splits $D=O\oplus O\theta$ with $\theta$ of
trace zero, so $\theta^2=c\in O$.  If $p^2$ divided $c$, then
$\theta/p$ would be integral and belong to $D$, contradicting this
basis.  Thus $v_p(c)$ is zero or one.

For any finite order $B\subset D$ containing $O$, half the trace
sends $B$ to $O\subset B$.  Its trace-zero part is therefore
$p^rO\theta$ for some $r\ge0$, and
$B=O\oplus p^rO\theta$.  Nonmaximality means $r\ge1$.
Consequently
\[
 B\simeq O[t]/(t^2-p^{2r}c),
\]
and the quadratic constant lies in $p^2O$.  This finite algebra is
local, since its quotient modulo $p$ is $\F_p[t]/(t^2)$.
Apply Theorem~\ref{thm:odd-quadratic-orders}.

The maximal order $D$ itself can be semilocal when $p$ splits in
the number field.  No nonsquare-residue claim about a unit $c$ is
needed here.
\end{proof}

\begin{example}\label{ex:cubic-quadratic-order}
Let $O=\Z_{(3)}$ or $\Z_3$ and take the nonmaximal ramified
quadratic order $A=O[t]/(t^2-27)=O[3\sqrt3]$.
The explicit choices $\alpha=-8$ and $\kappa=1/19$ satisfy
$\alpha^3=1-27/\kappa$.  They give
\[
 T=\begin{pmatrix}
 -8&0&19t\\
 t&0&-64\\
 0&1&8
 \end{pmatrix},\qquad T^3=I.
\]
The cyclic difference $[A^3,T]-[A[C_3]]$ has exact order two.
Its induction to $S_3$ is detected by $\mu_{u-1}$ after restriction
to $C_3$, since the other normalizer twist has eigenvalue $-1$.
\end{example}
